\section{Partition Equal Subset Sum}
\label{sec:dp-partition-equal-subset}

\problemstatement{Given an array $\textit{arr}$, determine whether it can
be partitioned into \textbf{two} subsets with \textbf{equal} sum.}

\begin{patternbox}
\emph{``Partition into two equal-sum subsets''} reduces directly to
Section~\ref{sec:dp-subset-sum-target}: if the whole array can be split
into two subsets of equal sum $S$, then the total sum must be $2S$ (so
an \textbf{odd} total sum makes this immediately impossible), and finding
one subset summing to exactly $S = \textit{totalSum}/2$ automatically
gives the other subset too (everything not in the first subset).
\end{patternbox}

\subsection*{Understanding the reduction}

No new recurrence needs deriving: compute $\textit{totalSum}$; if it is
odd, return \texttt{false} immediately. Otherwise, ask exactly
Section~\ref{sec:dp-subset-sum-target}'s question with target
$\textit{totalSum}/2$. If a subset summing to $\textit{totalSum}/2$
exists, the remaining elements necessarily sum to the same value
(since the two parts must add up to $\textit{totalSum} =
2 \times \textit{totalSum}/2$), giving a valid equal partition.

\subsection*{All Four Approaches}

Every one of the four stages is Section~\ref{sec:dp-subset-sum-target}'s
implementation, called with $k = \textit{totalSum}/2$. Only the
space-optimized version is shown here, to avoid repeating the identical
code four times.

\begin{lstlisting}
#include <bits/stdc++.h>
using namespace std;

bool subsetSumSpaceOptimized(vector<int>& arr, int k) {
    int n = arr.size();
    vector<bool> prev(k + 1, false);

    prev[0] = true;
    if (arr[0] <= k) prev[arr[0]] = true;

    for (int index = 1; index < n; index++) {
        vector<bool> current(k + 1, false);
        current[0] = true;
        for (int target = 1; target <= k; target++) {
            bool notPick = prev[target];
            bool pick = (arr[index] <= target) ? prev[target - arr[index]] : false;
            current[target] = notPick || pick;
        }
        prev = current;
    }
    return prev[k];
}

bool canPartition(vector<int>& arr) {
    int totalSum = accumulate(arr.begin(), arr.end(), 0);
    if (totalSum % 2 != 0) return false;

    return subsetSumSpaceOptimized(arr, totalSum / 2);
}

int main() {
    vector<int> arr = {1, 5, 11, 5};
    cout << (canPartition(arr) ? "true" : "false") << endl; // true
    return 0;
}
\end{lstlisting}

\subsection*{Dry run}

\dryrun{Take $\textit{arr} = [1,5,11,5]$.}

$\textit{totalSum} = 22$, even, so target $= 11$. Running
Section~\ref{sec:dp-subset-sum-target}'s subset-sum check for target
$11$: the subset $\{5,5,1\}$ sums to $11$ -- found, so the answer is
\texttt{true}. The complementary subset $\{11\}$ also sums to $11$,
confirming the equal partition $\{1,5,5\}$ and $\{11\}$.

\begin{complexitybox}
\textbf{Time Complexity.} $O(n \cdot \textit{totalSum})$, inherited
directly from Section~\ref{sec:dp-subset-sum-target}'s subset-sum DP
with $k = \textit{totalSum}/2$.
\textbf{Space Complexity.} $O(\textit{totalSum})$ for the rolling row
(space-optimized version).
\end{complexitybox}

\begin{takeawaybox}
Before deriving a new DP recurrence, always check whether the problem is
a thin wrapper around one already solved -- here, an odd-sum early exit
plus a single call to subset-sum with a derived target. This is the same
``reduce to an already-solved sub-problem'' instinct already exercised
for House Robber II's circular-array reduction back in
the DP 1D chapter.
\end{takeawaybox}
